% =====================================================
% 题型：圆台体积与轴截面选择题 (q_solid_frustum_volume.tex)
% =====================================================
% =====================================================
% 难度：★★ (中档)
% =====================================================
\newcommand{\qSolidFrustumVolumeStem}{已知圆台的上、下底面半径分别为 \(1\) 和 \(2\)，且圆台的母线与一底面所成的角的大小为 \(\dfrac{\pi}{4}\)，则圆台的体积为 \hfill （\quad）}

\newcommand{\qSolidFrustumVolumeOptions}{%
  \mcoptions[4]{\(\dfrac{4\pi}{3}\)}{\(\dfrac{7\pi}{3}\)}{\(\dfrac{5\pi}{2}\)}{\(3\pi\)}%
}

\newcommand{\qSolidFrustumVolumeAnswer}{B}

\newcommand{\qSolidFrustumVolumeFigure}[1][0.85]{%
  \begin{tikzpicture}[scale=#1, line join=round, line cap=round]
    % Lower ellipse
    \draw[thick] (2,0) arc[start angle=0, end angle=-180, x radius=2, y radius=0.6];
    \draw[thick, dashed] (-2,0) arc[start angle=180, end angle=0, x radius=2, y radius=0.6];
    % Upper ellipse
    \draw[thick] (0,1.2) ellipse[x radius=1, y radius=0.3];
    % Outer boundary lines (generators)
    \draw[thick] (-2,0) -- (-1,1.2);
    \draw[thick] (2,0) -- (1,1.2);
    
    % Centers
    \coordinate (O1) at (0,0);
    \coordinate (O2) at (0,1.2);
    \coordinate (A) at (2,0);
    \coordinate (B) at (1,1.2);
    \coordinate (C) at (1,0);
    
    \ifteacher
      % Auxiliary lines for teacher
      \draw[thick, dashed, gray] (O1) -- (O2) node[midway, left] {\small \(h\)};
      \draw[thick, dashed, gray] (O2) -- (B);
      \draw[thick, dashed, gray] (B) -- (C) node[midway, right] {\small \(h\)};
      \draw[thick, dashed, gray] (O1) -- (A) node[midway, below] {\small \(R=2\)};
      \draw[thick, dashed, gray] (C) -- (A) node[midway, below] {\small \(1\)};
      \fill[gray] (C) circle (1.2pt) node[below] {\small \(C\)};
      \draw[gray] (C) ++(0,0.15) -- ++(0.15,0) -- ++(0,-0.15);
      \draw[->, gray] (A) +(-0.4,0) arc[start angle=180, end angle=135, radius=0.4];
      \node[gray] at (1.4,0.2) {\small \(45^\circ\)};
    \else
      % Simple centers and radii for student
      \draw[thick, dashed] (O1) -- (A);
      \draw[thick, dashed] (O2) -- (B);
    \fi
    
    \fill (O1) circle (1.2pt) node[below] {\small \(O_1\)};
    \fill (O2) circle (1.2pt) node[above] {\small \(O_2\)};
    \fill (A) circle (1.2pt) node[right] {\small \(A\)};
    \fill (B) circle (1.2pt) node[above right] {\small \(B\)};
  \end{tikzpicture}%
}

\newcommand{\qSolidFrustumVolumeSolution}{%
  圆台体积公式为：
  \[
    V = \frac{1}{3}\pi h (R^2 + Rr + r^2)
  \]
  其中下底半径 \(R = 2\)，上底半径 \(r = 1\)。我们只需计算圆台的高 \(h\)。
  
  圆台的轴截面是一个等腰梯形。做上底面圆周上一点对下底面的投影，
  在形成的直角三角形中：
  - 直角边之一为高 \(h\)，
  - 另一条直角边为下底半径与上底半径之差，即 \(R - r = 2 - 1 = 1\)，
  - 母线与底面所成的角为 \(\dfrac{\pi}{4}\)。
  
  因此有：
  \[
    h = (R - r) \cdot \tan\frac{\pi}{4} = 1 \cdot 1 = 1.
  \]
  
  将 \(h = 1\), \(R = 2\), \(r = 1\) 代入圆台体积公式：
  \[
    V = \frac{1}{3}\pi \cdot 1 \cdot (2^2 + 2 \cdot 1 + 1^2) = \frac{1}{3}\pi (4 + 2 + 1) = \frac{7\pi}{3}.
  \]
  故选 B。%
}
